\begin{figure}
  \centering
  \begin{tikzpicture}[scale=0.15]
    \draw[fill=MaterialColor,very thick,draw=MaterialColor] (0,0) rectangle (6,6);
    \draw[fill=MaterialColor,very thick,draw=MaterialColor] (9,0) rectangle (14,9);
    \draw[fill=MaterialColor,very thick,draw=MaterialColor] (3,9) rectangle (14,18);
    \draw[fill=VoidColor,very thick,draw=VoidColor] (0,6) rectangle (3,18);
    \draw[fill=VoidColor,very thick,draw=VoidColor] (3,6) rectangle (9,9);
    \draw[fill=VoidColor,very thick,draw=VoidColor] (6,0) rectangle (9,9);
    \draw[fill=SourceColor,very thick,draw=SourceColor] (0,15) rectangle (3,18);

    \draw[fill=none, very thick] (0,0) rectangle (14,18);

    \node[draw=none, fill=none, rotate=90] at (-1,9) {Reflective Boundary};
    \node[draw=none, fill=none, rotate=0] at (7,19) {Reflective Boundary};
    \node[draw=none, fill=none, rotate=-90] at (15,9) {Vacuum Boundary};
    \node[draw=none, fill=none, rotate=0] at (7,-1) {Vacuum Boundary};

    \draw[fill=SourceColor,very thick,align=left] (19,11.5) rectangle ++(2.25,2.25);
    \node[anchor=west] at (21.7,12.5) {= Source};

    \draw[fill=VoidColor,very thick,align=left] (19,7.0) rectangle ++(2.25,2.25);
    \node[anchor=west] at (21.7,8.25) {= Void};

    \draw[fill=MaterialColor,very thick,align=left] (19,2.5) rectangle ++(2.25,2.25);
    \node[anchor=west] at (21.7,3.5) {= Absorbing/Scattering Material};
  \end{tikzpicture}
  \caption{\footnotesize Schematic of \gls{arp3}.}
\end{figure}